%!TEX program = uplatex
\RequirePackage{plautopatch}
\documentclass[uplatex,dvipdfmx]{jlreq}
\usepackage{amsmath,amssymb}

\pagestyle{empty}
\begin{document}

\noindent
{\large\bfseries 数論的トロイダルコンパクト化}

\medskip
\hspace*{\fill}{\large\bfseries 藤原 一宏（名大・多元数理）}

\bigskip\bigskip

有界対称領域の算術商をみると、数学者は胸がときめく。
この講演では数論の人間のときめきを語ってみたい。
まず、予備知識をいくつか。
そのような多様体にはある標準的な仕方で代数体上の代数多様体の構造が入ることが、
志村五郎氏（プリンストン大）の一連の仕事、及びそれを補完する多くの人たちの仕事により
示されている（それ故に志村多様体といわれる）。
この事実は現在の数論において極めて重要な役割をはたしている。

まず、これが虚数乗法論の一般化（の一つの方向）であること、
つまりモデュラー関数の特殊値がどのようなアーベル拡大を生成するかに解答を与えている
（Kronecker の青春の夢、Hilbert の第12問題）。
また、そのエタールコホモロジー群上にはガロア群とアデール化された代数群が共に作用し、
ガロア表現と保型表現との間の非可換類体論を実現すると信じられている
（コホモロジーの研究のためにもコンパクト化が重要なのである）。

\medskip

志村多様体自身の話題に戻ろう。
志村氏の仕事にもアーベル多様体という形で現れているが、
Deligne は Grothendieck のモチーフの哲学
（モチーフという数学的実在を信じる宗教のようなもの）と志村氏の仕事を結びつけた。
夢を語るとすれば、志村多様体は（条件を付けた）モチーフのモデュライ空間なのである。
（有界対称領域自身は Hodge 構造の分類空間と解釈される。「モチーフは Hodge 構造を生む」のである。）

このようなモチーフ論的視点に立つと、
トロイダルコンパクト化の構成はモチーフの退化理論の帰結としてとらえられることになる
（残念ながら、整数環 $\mathbf{Z}$ 上満足のいく形でのモチーフ理論は知られていないのだが）。
ここでは特殊なモチーフであるアーベル多様体とその退化である $1$-motive を通して
PEL 型の志村多様体のトロイダルコンパクト化が $\mathbf{Z}$ 上自然に構成できることを述べたい
（Siegel modular 多様体の場合は Chai 及び Faltings により、PEL 型の時は筆者による）。
しかしながら、あまり一般的には行わず、
楕円モデュラー、ヒルベルトモデュラーなどのより親しみやすい例から入りたいと思う。

\end{document}